
\include{./def/yf-formatting}
\include{./def/yf-def}

\usepackage{wasysym}
%\usepackage{times}

\title{}
\author{}
\date{}


%\def\A{\mathcal{A}}
%\def\B{\mathcal{B}}
\def\C{\mathcal{C}}
%\def\D{\mathcal{D}}
\def\E{\mathcal{E}}
\def\FF{\EuScript{F}}
\def\G{\mathcal{G}}
\def\GG{\EuScript{G}}
\def\H{\mathcal{H}}
\def\I{\mathcal{I}}
\def\J{\mathcal{J}}
\def\L{\mathcal{L}}
\def\Q{\mathcal{Q}}
\def\S{\mathcal{S}}
\def\T{\mathcal{T}}
\def\V{\mathcal{V}}
%\def\W{\EuScript{W}}
\def\W{\mathscr{W}}
\def\X{\mathcal{X}}
\def\Y{\mathcal{Y}}
\def\Z{\mathcal{Z}}
\def\att{\mit{diffatt}}
\def\boldatt{\mathbf{att}}
\def\ext{\mathit{ext}}
\def\join{\mathit{join}}
\def\schema{\mathit{schema}}


\def\extraspacing{\vspace{4mm} \noindent}
\def\vgap{\vspace{4mm}}
\def\vslit{\vspace{2mm}}

\begin{document}

\boxminipg{\linewidth}{
    \begin{center}
        \vspace{3mm}
    {\Large
    (CE7456 Lecture Notes 10) Natural Joins} \\[2mm]
    {\large Yufei Tao}
    \vspace{3mm}
    \end{center}
}

Today, we discuss the {\em natural join} problem, which lies at the core of database theory and generalizes the subgraph enumeration problem in the previous lecture.

\vgap

Let $\boldatt$ be a finite set, where each element is called an {\em attribute}. For a non-empty set $\X \subseteq \boldatt$ of attributes, a {\em tuple} over $\X$ is a function $\bm{u}: \X \rightarrow \intdom$ (where $\intdom$ represents the set of integers); for each $X \in \X$, we refer to $\bm{u}(X)$ as the {\em value} of $\bm{u}$ on $X$. For any non-empty subset $\Y \subseteq \X$, we define $\bm{u}[\Y]$ --- called the {\em projection} of $\bm{u}$ on $\Y$ --- as the tuple $\bm{v}$ over $\Y$ satisfying $\bm{v}(Y) = \bm{u}(Y)$ for every attribute $Y \in \Y$. A {\em relation} $R$ is a set of tuples over the same set $Z$ of attributes; we refer to $Z$ as the {\em schema} of $R$ and represent it as $\schema(R)$.

\vgap

We define a {\em join} as a set $\Q$ of relations. Let $\schema(\Q) = \bigcup_{R \in Q} \schema(R)$. The result of $\Q$ is a relation over $\schema(\Q)$ formalized as:
\myeqn{
    \join(\Q) &=& \{\textrm{tuple $\bm{u}$ over $\schema(\Q)$} \mid \forall R \in \Q: \bm{u}[\schema(R)] \in R\}. \label{eqn:join-rslt}
}
Our discussion in this lecture assumes that $\schema(\Q)$ has a constant size.

\extraspacing {\bf Math Conventions.} Define $0^0 = 0$. For an integer $x \ge 1$, the notation $[x]$ denotes the set $\{1, 2, ..., x\}$.

%Every logarithm $\log(\cdot)$ has base $2$.

%The notation $\tO(.)$ hides a factor logarithmic to $\sum_{R \in \Q} |R|$.

\section{Schema Graphs and Fractional Edge Covers} \label{sec:schema}

Given a join $\Q$, we define its {\em schema graph} as the hypergraph $\H = (\V, \E)$ where
\myeqn{
    \V &=& \schema(\Q) \nn \\
    \E &=& \set{\schema(R) \mid R \in \Q}. \nn
}
Note that there is one-one correspondence between the relations in $\Q$ and the hyperedges in $\E$. For each hyperedge $e \in \E$, let $R_e$ represent the corresponding relation $R \in \Q$.

\vgap

Treating $\E$ as a ``space'', next we define vectors over it. Every such vector $\W$ has a real-valued component $\W_e$ on each hyperedge $e \in \E$. The vector $\W$ is said to be a {\em fractional edge cover} of $\H$ if
\myitems{
    \item $0 \le \W_e \le 1$ for every $e \in \E$, and
    \item $\sum_{e \in \E : X \in e} \W_e \ge 1$ for every $X \in \V$.
}

\section{AGM Bound} \label{sec:agm}

Let $\Q$ be a join with schema graph $\H = (\V, \E)$. Every fractional edge cover $\W$ of $\H$ is associated with an {\em AGM value} defined as
\myeqn{
    \mit{AGM}_\W &=& \prod_{e \in \E} |R_e|^{\W_e}. \label{eqn:agm}
}

\vgap

\begin{theorem} \label{thm:agm}
   For any fractional edge cover $\W$ of $\H$, it holds that $|\join(\Q)| \le \mit{AGM}_\W$.
\end{theorem}

\vgap

%An {\em optimal fractional edge cover} of $\H$ is a fractional edge cover $\W^*$ minimizing \eqref{eqn:agm}.
The {\em AGM bound} of $\Q$ is defined as the AGM value of $\W^*$.
The AGM bound is tight:

\vgap

\begin{theorem} \label{thm:agm-tight}
   Fix an arbitrary hypergraph $\H = (\V, \E)$ and an arbitrary integer vector $N$ over $\E$ such that $N_e \ge 1$ for each $e \in \E$. There exists a join $\Q$ satisfying
   \myitems{
        \item the schema graph of $\Q$ is $\H$;
        \item for each $e \in \E$, $|R_e| = N_e$;
        \item $|\join(\Q)|$ is at least the AGM bound of $\Q$ up to a constant factor.
   }
\end{theorem}

\vgap

We prove Theorem~\ref{thm:agm} in the next section and Theorem~\ref{thm:agm-tight} in the next lecture.

\section{Proof of Theorem~\ref{thm:agm}} \label{sec:agm-proof}

Given a tuple $\bm{u}$ over a subset $U \subseteq \E$ (recall that $\E = \schema(\Q)$) and a relation $R \in \Q$, we define a {\em residual relation} as:
\myeqn{
    R(\bm{u}) &=&
    \set{
        \bm{v} \in R \mid \forall X \in U \cap \schema(R): \bm{v}(X) = \bm{u}(X)
    }. %\nn
    \label{eqn:alg-proof:residual-rel}
}
Intuitively, $R(\bm{u})$ includes those tuples of $R$ that are ``constrained'' by $\bm{u}$, in the sense that they must take the same value as $\bm{u}$ on every attribute common to $\schema(R)$ and $U$. Note that if $U$ and $\schema(R)$ are disjoint, then $R(\bm{u}) = R$.

\vgap

Given a tuple $\bm{u}$ over a subset $U \subseteq \E$, we define a {\em residual join} as follows:
\myeqn{
    \Q(\bm{u}) &=&
    \set{
        R_e(\bm{u}) \mid e \in \E
    }. \label{eqn:alg-proof:residual-join}
}

\vslit

\begin{lemma} \label{lmm:agm-proof:residual}
    Fix an arbitrary order of the attributes in $\V$: $X_1, X_2, ..., X_d$ where $d = |\V|$. Let $t \in [0, d]$ be an integer and $\bm{u}$ be a tuple over $U = \set{X_1, ..., X_t}$ such that $R(\bm{u}) \ne \emptyset$ for every $R \in \Q$. We have
    \myeqn{
        |\join(\Q(\bm{u}))|
        &\le&
        \prod_{e \in \E} |R_e(\bm{u})|^{\W_e}. \label{eqn:agm-proof:residual-goal}
    }
\end{lemma}

\vgap

Specially, if $t = 0$, then $U = \emptyset$ and $\Q = \Q(\bm{u})$ (because $R(\bm{u}) = R$ for all $R \in \Q$). Hence, Theorem~\ref{thm:agm} is a corollary of the above lemma with $U = \emptyset$. Note that the theorem is still correct even if $R(\bm{u}) = \emptyset$ for some $R \in \Q$. We choose to state the theorem the way it is because this will bring us convenience when we prove Lemma~\ref{lmm:alg:residual} later.

\extraspacing {\bf Proof of Lemma~\ref{lmm:agm-proof:residual}.} In general, given a relation $R$ and an attribute $X$ in its schema, we define
\myeqn{
    \pi_X (R) &=& \set{\bm{v}(X) \mid \bm{v} \in R}. \label{eqn:agm-proof:proj}
}
This is the {\em projection} operator in relational algebra.

\vgap

We prove the lemma by induction. As the base case, consider $t = d$. For each $R \in \Q$, the residual relation $R(\bm{u})$ has exactly one tuple. Thus, $|\join(\Q(\bm{u}))| \le 1$, and \eqref{eqn:agm-proof:residual-goal} holds.

\vgap

Inductively, assuming that the lemma holds for $t = i$ (where $i \in [d]$), next we prove its correctness for $t = i - 1$. Because $\W$ is a fractional edge cover of $\H = (\V, \E)$, we know
\myeqn{
    \sum_{e \in \E: X_i \in e} &\ge& 1.
    \label{eqn:agm-proof:residual-goal:3}
}
Define:
\myeqn{
    \I = \bigcap_{e \in \E : X_i \in e} \pi_{X_i}(R_e(\bm{u})). \label{eqn:agm-proof:intersection}
}
For each value $x \in \I$, we construct a tuple $\bm{v}$ over $\set{X_1, ..., X_i}$ by ``extending $\bm{u}$ with $X_i = x$'', i.e., $\bm{v}(X_j) = \bm{u}(X_j)$ for all $j \in [i - 1]$ and $\bm{v}(X_i) = x$. This produces a set of ``extended'' tuples:
\myeqn{
    S_\ext &=&
    \set{
        \text{tuple $\bm{v}$ extended from $\bm{u}$ and $x$} \mid
        \text{$x \in \I$}
    }. \label{eqn:agm-proof:S_ext}
}

\vslit

Consider the residual join $\Q(\bm{v})$ where $\bm{v} \in S_\ext$. We argue that no relation of $\Q(\bm{v})$ is empty. Indeed, for each $e \in \E$:
\myitems{
    \item if $X_i \notin e$, then $R_e(\bm{v}) = R_e(\bm{u})$, which is not empty;

    \item otherwise, since $\bm{v}(X_i)$ belongs to the set $\I$ in \eqref{eqn:agm-proof:intersection}, we can assert $R_e(\bm{v}) \ne \emptyset$.
}

\vslit

We can now derive
\myeqn{
    \join(\Q(\bm{u}))
    &=&
    \bigcup_{\bm{v} \in S_\ext} \join(\Q(\bm{v})) \nn \\
    \Rightarrow
    |\join(\Q(\bm{u}))|
    &=&
    \sum_{\bm{v} \in S_\ext} |\join(\Q(\bm{v}))| \nn \\
    \explain{by inductive assumption}
    &\le&
    \sum_{\bm{v} \in S_\ext}
    \prod_{e \in \E}
    |R_e(\bm{v})|^{\W_e} \nn \\
    &\le&
    \sum_{\bm{v} \in S_\ext}
    \Big(
    \prod_{e \in \E : X_i \in e}
    |R_e(\bm{v})|^{\W_e}
    \Big)
    \Big(
    \prod_{e \in \E : X_i \notin e}
    |R_e(\bm{v})|^{\W_e}
    \Big) \nn \\
    &=&
    \sum_{\bm{v} \in S_\ext}
    \Big(
    \prod_{e \in \E : X_i \in e}
    |R_e(\bm{v})|^{\W_e}
    \Big)
    \Big(
    \prod_{e \in \E : X_i \notin e}
    |R_e(\bm{u})|^{\W_e}
    \Big) \nn \\
    &=&
    \Big(
    \prod_{e \in \E : X_i \notin e}
    |R_e(\bm{u})|^{\W_e}
    \Big)
    \sum_{\bm{v} \in S_\ext}
    \Big(
    \prod_{e \in \E : X_i \in e}
    |R_e(\bm{v})|^{\W_e}
    \Big) \nn \\
    &\le&
    \Big(
    \prod_{e \in \E : X_i \notin e}
    |R_e(\bm{u})|^{\W_e}
    \Big)
    \prod_{e \in \E: X_i \in e}
    \Big(
        \sum_{\bm{v} \in S_\ext} |R_e(\bm{v})|
    \Big)^{\W_e} \nn \\
    &&
    \explain{by \eqref{eqn:agm-proof:residual-goal:3} and Friedgut's inequality \eqref{eqn:friedgut} in Appendix \ref{app:friedgut}}
    \nn \\[2mm]
    &\le&
    \Big(
    \prod_{e \in \E : X_i \notin e}
    |R_e(\bm{u})|^{\W_e}
    \Big)
    \prod_{e \in \E: X_i \in e}
        |R_e(\bm{u})|^{\W_e} \nn \\
    &=&
    \prod_{e \in \E}
    |R_e(\bm{u})|^{\W_e}. \label{eqn:agm-proof:residual-goal:4}
}
This completes the inductive step.


\section{A Join Algorithm} \label{sec:alg}

Let $\Q$ be a join with schema graph $\H = (\V, \E)$. Remarkably, the proof of the AGM bound inspires an algorithm for computing $\join(\Q)$. Set $d = |\V|$ and fix an arbitrary order of the attributes in $\V$: $X_1, X_2, ..., X_d$. We make two assumptions:
\myitems{
    \item {{\bf A1:}} Given a tuple $\bm{u}$ over a set $U = \set{X_1, ..., X_t}$ where $t \in [0, d]$ and a relation $R \in \Q$, we can compute the residual relation $R(\bm{u})$ --- defined in \eqref{eqn:alg-proof:residual-rel} --- in $O(1 + |R(\bm{u})|)$ time and $|R(\bm{u})|$ in $O(1)$ time.

    \item {{\bf A2:}}  Given a tuple $\bm{u}$ over a set $U = \set{X_1, ..., X_t}$ where $t \in [0, d-1]$, a relation $R \in \Q$, we can test in $O(1)$ time whether $R(\bm{u}) = \emptyset$.
}
These assumptions can be removed by slowing down the running time by at most a logarithmic factor (details left to you).

\vgap

We aim to solve the following {\em residual join computation problem}:

\vslit

\minipg{0.9\linewidth}{
    Let $t \in [0, d]$ be an integer. Given a tuple $\bm{u}$ over $U = \set{X_1, ..., X_t}$ such that $R(\bm{u}) \ne \emptyset$ for each $R \in \Q$, compute the result of the residual join $\Q(\bm{u})$ defined in \eqref{eqn:alg-proof:residual-join}.
}

\vslit

\noindent Computing $\join(\Q)$ is equivalent to solving the above with $t = 0$ (and hence $U = \emptyset$).

\extraspacing {\bf Algorithm.} The base case $t = d$ can be trivially settled in $O(1)$ time (assumption {\bf A1}). Inductively, suppose that we can solve the problem for $t = i$ for some $i \in [d]$. We solve the problem for $t = i-1$ as follows:
\myenums{
    \item Find the set $\I$ in \eqref{eqn:agm-proof:intersection} and construct the set $S_\ext$ in \eqref{eqn:agm-proof:S_ext}.

    \item For each $\bm{v} \in S_\ext$, output $\join(\Q(\bm{v}))$.
}

\extraspacing {\bf Analysis.} We prove:

\begin{lemma} \label{lmm:alg:residual}
    Fix an arbitrary fractional edge cover $\W$ of $\H$.
    The residual join computation problem can be solved in time
    \myeqn{
        O\Big(
        \prod_{e \in \E} |R_e(\bm{u})|^{\W_e}
        \Big) \label{eqn:alg:residual-time}
    }
    under assumptions {\bf A1} and {\bf A2}.
\end{lemma}

\begin{proof}
   The base case $t = d$ is straightforward and omitted. Next, assuming that the claim holds for $t = i \in [d]$, we prove the claim for $t = i - 1$.

   \vgap

   In Step 1 of the algorithm, the set $\I$ in \eqref{eqn:agm-proof:intersection} (and, hence, $S_\ext$ in \eqref{eqn:agm-proof:S_ext}) can be computed in
   \myeqn{
        O\Big(
            \min_{e \in \E: X_i \in e} |R_e(\bm{u})|
        \Big) \label{eqn:alg:residual:1}
   }
   time as follows. First, spend $O(1)$ time to find the edge $e^* = \argmin_{e \in \E: X_i \in e} |R_e(\bm{u})|$ (this is possible due to assumption {\bf A1}). Then, retrieve the tuples in $R_{e^*}(\bm{u})$, which takes $O(|R_{e^*}(\bm{u})|)$ time by assumption {\bf A1}. For each tuple in $R_{e^*}(\bm{u})$, we test in $O(1)$ time whether the tuple's value on $X_i$ belongs to $\I$ (this is possible due to assumption {\bf A2}). The complexity in \eqref{eqn:alg:residual:1} now follows.

   \vgap

   Next, we show that \eqref{eqn:alg:residual:1} is bounded by \eqref{eqn:alg:residual-time}. Because $\W$ is a fractional edge cover of $\H = (\V, \E)$, it holds that
    \myeqn{
        \sum_{e \in \E : X_i \in e} \W_e &\ge& 1.
        \label{eqn:alg:residual:2}
    }
    Hence:
    \myeqn{
        \min_{e \in \E: X_i \in e} |R_e(\bm{u})|
        &\le&
        \Big(
            \min_{e \in \E : X_i \in e} |R_e(\bm{u})|
        \Big)^{\sum_{e \in \E : X_i \in e} \W_e} \nn \\
        &=&
        \prod_{e \in \E : X_i \in e}
        \Big(
            \min_{e \in \E : X_i \in e} |R_e(\bm{u})|
        \Big)^{\W_e} \nn \\
        &\le&
        \prod_{e \in \E : X_i \in e} |R_e(\bm{u})|^{\W_e}. \nn
    }

    \vslit

    Let us now turn to Step 2 of the algorithm. By the inductive assumption, for each $\bm{v} \in S_\ext$, the result of $\Q(\bm{v})$ can be computed in $O(\prod_{e \in \E} |R_e(\bm{v})|^{\W_e})$ time. Therefore, the total cost of Step 2 is
    \myeqn{
        O\Big(
            \sum_{\bm{v} \in S_\ext}
            \prod_{e \in \E} |R_e(\bm{v})|^{\W_e}
        \Big) \nn
    }
    which, by the derivation in \eqref{eqn:agm-proof:residual-goal:4}, is bounded by \eqref{eqn:alg:residual-time}.
\end{proof}

\begin{corollary}
   Under assumptions {\bf A1} and {\bf A2}, we can compute the result of a join $\Q$ in time that is asymptotically bounded by the AGM bound of $\Q$.
\end{corollary}

\section{Remark}

The AGM bound was first established in \cite{agm13}. The join algorithm presented is a special case of an algorithm given in \cite{nrr13}.


\bibliographystyle{plainurl}% the mandatory bibstyle
\bibliography{ref}

\appendix

\section{Appendix: Friedgut's Inequality} \label{app:friedgut}

Fix arbitrary integers $p$ and $q$ at least 1. Take a set of non-negative real values $\{a_{i,j} \mid i \in [p], j \in [q]\}$. In addition, take another set of non-negative real values $\{b_i$ $\mid$ $i \in [q]\}$ satisfying $\sum_{i=1}^{q} b_i \geq 1$.  Friedgut's inequality \cite{f04} states that
    \myeqn{
        \sum_{i=1}^{p}	\prod_{j=1}^{q} a_{i, j}^{b_j}
        &\leq&
        \prod_{j=1}^{q} \left(\sum_{i=1}^{p} a_{i, j}\right)^{b_j}.
        \label{eqn:friedgut}
    }

\end{document}
